\documentclass[aspectratio=169, 11pt]{beamer}

% ============================================================
% Pacchetti
% ============================================================
\usepackage[utf8]{inputenc}
\usepackage[italian]{babel}
\usepackage[T1]{fontenc}
\usepackage{lmodern}
\usepackage{eurosym}
\usepackage{tikz}

% ============================================================
% Tema e colori (palette dell'app)
% ============================================================
\usetheme{Madrid}
\usecolortheme{default}

\definecolor{appblue}{RGB}{21,101,192}
\definecolor{appdark}{RGB}{0,60,143}
\definecolor{appaccent}{RGB}{255,112,67}
\definecolor{appgreen}{RGB}{67,160,71}
\definecolor{applight}{RGB}{233,242,251}

% Applica la palette ai vari elementi del tema
\setbeamercolor{palette primary}{bg=appblue, fg=white}
\setbeamercolor{palette secondary}{bg=appdark, fg=white}
\setbeamercolor{palette tertiary}{bg=appdark, fg=white}
\setbeamercolor{frametitle}{bg=appblue, fg=white}
\setbeamercolor{title}{fg=appdark}
\setbeamercolor{structure}{fg=appblue}
\setbeamercolor{block title}{bg=appblue, fg=white}
\setbeamercolor{block body}{bg=applight, fg=black}
\setbeamercolor{item}{fg=appblue}
\setbeamertemplate{navigation symbols}{} % nasconde i simboli di navigazione
\setbeamertemplate{frametitle}[default][left]

% Comando di comodo per una "card" colorata
\newcommand{\card}[2]{%
  \begin{beamercolorbox}[wd=\linewidth, sep=6pt, rounded=true, shadow=true]{block body}%
    \textbf{\textcolor{appblue}{#1}}\\[2pt]\small #2%
  \end{beamercolorbox}%
}

% ============================================================
% Metadati
% ============================================================
\title[Travel Organizer Offline]{Travel Organizer Offline}
\subtitle{App mobile in Flutter per organizzare i viaggi, anche offline}
\author[Gruppo]{Gruppo \\ \small Componente 1 \and Componente 2 \and Componente 3}
\institute[]{Mobile Programming -- A.A. 2025/2026}
\date{Giugno 2026}

\begin{document}

% ---------- Frontespizio ----------
\begin{frame}[plain]
  \titlepage
  \begin{center}
    \begin{tikzpicture}
      \fill[appblue, rounded corners=8pt] (0,0) rectangle (1.1,1.1);
      \fill[white, rounded corners=2pt] (0.3,0.28) rectangle (0.8,0.72);
      \fill[appblue] (0.3,0.5) rectangle (0.8,0.56);
      \draw[white, line width=1.2pt, rounded corners=2pt] (0.42,0.72) rectangle (0.68,0.86);
    \end{tikzpicture}
  \end{center}
\end{frame}

% ---------- Indice ----------
\begin{frame}{Indice}
  \tableofcontents
\end{frame}

% ============================================================
\section{L'idea}
\begin{frame}{L'idea dell'app}
  \begin{block}{Obiettivo}
    Aiutare l'utente a \textbf{organizzare i propri viaggi} -- itinerari, tappe,
    attività, checklist e spese -- mantenendo tutto disponibile
    \textbf{anche senza connessione di rete}.
  \end{block}
  \vfill
  \begin{columns}[T]
    \begin{column}{0.5\textwidth}
      \card{Il problema}{Le info di un viaggio sono sparse tra note, chat e siti
      e richiedono una connessione per essere consultate.}
    \end{column}
    \begin{column}{0.5\textwidth}
      \card{La soluzione}{Un'unica app, \textbf{100\% offline}, senza backend né
      API remote: i dati vivono sul dispositivo.}
    \end{column}
  \end{columns}
  \vfill
  \footnotesize \textbf{Utenti:} viaggiatori, studenti, famiglie -- chiunque
  voglia tenere sotto controllo programma, bagaglio e budget.
\end{frame}

% ============================================================
\section{Funzionalità principali}
\begin{frame}{Funzionalità principali}
  \begin{columns}[T]
    \begin{column}{0.5\textwidth}
      \begin{itemize}
        \item \textbf{Viaggi} -- CRUD, stato, budget, tag, modalità di viaggio
        \item \textbf{Tappe / giornate} -- riordinabili con drag \& drop
        \item \textbf{Attività} -- categoria, orario, luogo, costo, stato
        \item \textbf{Checklist} -- documenti e cose da fare, con avanzamento
      \end{itemize}
    \end{column}
    \begin{column}{0.5\textwidth}
      \begin{itemize}
        \item \textbf{Spese} -- previste/effettive, budget
        \item \textbf{Ricerca e filtri} -- su ogni entità + ricerca globale
        \item \textbf{Statistiche} -- riepiloghi e grafici
        \item \textbf{Calendario} -- vista mensile con marker
      \end{itemize}
    \end{column}
  \end{columns}
  \vfill
  \begin{beamercolorbox}[wd=\linewidth, sep=5pt, rounded=true]{block title}
    \centering Tutti i dati gestiti con \textbf{inserimento, modifica,
    visualizzazione ed eliminazione} (CRUD completo).
  \end{beamercolorbox}
\end{frame}

% ============================================================
\section{Schermate principali}
\begin{frame}{Le schermate più importanti}
  \textbf{Navigazione a 4 sezioni} (Bottom Navigation Bar):
  \vspace{0.3cm}
  \begin{columns}[T]
    \begin{column}{0.25\textwidth}
      \card{Viaggi}{Elenco con ricerca, filtri e azioni rapide.}
    \end{column}
    \begin{column}{0.25\textwidth}
      \card{Calendario}{Mese con marker per viaggi e attività.}
    \end{column}
    \begin{column}{0.25\textwidth}
      \card{Statistiche}{Totali, grafici e salute del budget.}
    \end{column}
    \begin{column}{0.25\textwidth}
      \card{Emergenza}{Dati medici e numeri utili.}
    \end{column}
  \end{columns}
  \vspace{0.4cm}
  \textbf{Dettaglio viaggio} -- 6 schede:
  \begin{center}
    \small Tappe \, $\cdot$ \, Attività \, $\cdot$ \, Checklist \, $\cdot$ \,
    Spese \, $\cdot$ \, Timeline \, $\cdot$ \, Valigia
  \end{center}
\end{frame}

% ============================================================
\section{Scelte progettuali}
\begin{frame}{Scelte progettuali}
  \begin{block}{Architettura a livelli}
    \centering
    UI (Schermate) $\rightarrow$ Provider (Stato) $\rightarrow$ Repository
    $\rightarrow$ Database SQLite
  \end{block}
  \vspace{0.2cm}
  \begin{columns}[T]
    \begin{column}{0.5\textwidth}
      \textbf{Pattern mobile}
      \begin{itemize}\small
        \item Stato con \texttt{provider} (ChangeNotifier)
        \item Navigazione \texttt{Navigator} + passaggio dati
        \item \texttt{IndexedStack} per i tab
        \item Material 3, empty state, validazioni
      \end{itemize}
    \end{column}
    \begin{column}{0.5\textwidth}
      \textbf{Persistenza}
      \begin{itemize}\small
        \item \texttt{sqflite} (locale, offline)
        \item Chiavi esterne + eliminazione a cascata
        \item Schema \textbf{versionato} con migrazioni
        \item Repository per le query SQL
      \end{itemize}
    \end{column}
  \end{columns}
\end{frame}

% ============================================================
\section{Feature avanzate}
\begin{frame}{Feature avanzate (1/2)}
  \begin{columns}[T]
    \begin{column}{0.5\textwidth}
      \card{Packing list intelligente}{Tag del viaggio (Mare, Estero…)
      $\Rightarrow$ oggetti suggeriti, \textbf{organizzati in sacchetti} con
      avanzamento per sacchetto.}
      \vspace{4pt}
      \card{Budget ``gamification''}{Barra verde$\to$giallo$\to$rosso +
      \textbf{ritmo di spesa}: ricalcola il budget giornaliero e avvisa sugli
      sforamenti.}
    \end{column}
    \begin{column}{0.5\textwidth}
      \card{Timeline con tempi liberi}{Attività come blocchi colorati per
      categoria; badge \emph{``Hai 3h di tempo libero''}.}
      \vspace{4pt}
      \card{Esporta / Importa viaggio}{Viaggio in JSON condiviso col menu
      nativo e reimportabile su un altro dispositivo.}
    \end{column}
  \end{columns}
\end{frame}

\begin{frame}{Feature avanzate (2/2)}
  \begin{columns}[T]
    \begin{column}{0.5\textwidth}
      \card{Diario di viaggio con foto}{Per ogni attività: un ricordo + una foto
      dalla galleria/fotocamera (\texttt{image\_picker}).}
      \vspace{4pt}
      \card{Notifiche locali}{Promemoria \textbf{30 min prima} di un'attività,
      anche ad app chiusa (API native, nessun server).}
      \vspace{4pt}
      \card{Calendario interattivo}{Marker su partenze, ritorni e attività.}
    \end{column}
    \begin{column}{0.5\textwidth}
      \card{Sezione Emergenza (ICE)}{Dati medici locali + numeri di emergenza
      \textbf{dedotti dal paese} di destinazione (112, 911, 110/119…).}
      \vspace{4pt}
      \card{Ricerca globale ``Spotlight''}{Cerca in tempo reale tra viaggi,
      tappe, attività e spese, con risultati raggruppati.}
      \vspace{4pt}
      \card{Duplicazione}{Di viaggi, tappe e checklist.}
    \end{column}
  \end{columns}
\end{frame}

% ============================================================
\section{Dati e integrità}
\begin{frame}{Integrità e coerenza dei dati}
  \begin{itemize}
    \item \textbf{Niente viaggi sovrapposti}: bloccate le date che si accavallano.
    \item \textbf{Eliminazione a cascata}: cancellando un viaggio si rimuovono
    tappe, attività, checklist e spese; le statistiche si aggiornano subito.
    \item \textbf{Validazioni}: date coerenti e campi obbligatori.
    \item \textbf{Test automatici}: avvio app, integrità date, esporta/importa,
    dati di esempio.
  \end{itemize}
  \vfill
  \begin{beamercolorbox}[wd=\linewidth, sep=5pt, rounded=true]{block title}
    \centering 4 viaggi di esempio precaricati: Tokyo, Barcellona, Roma
    (conclusi) e \textbf{Parigi} (itinerario completo, 27 ott -- 1 nov 2026).
  \end{beamercolorbox}
\end{frame}

% ============================================================
\section{Difficoltà e demo}
\begin{frame}{Difficoltà incontrate}
  \begin{itemize}
    \item \textbf{Integrazioni native multipiattaforma} (notifiche, fotocamera):
    gestite con ``guardie'' che le rendono inoffensive su Web e nei test.
    \item \textbf{Fusi orari} per le notifiche programmate: risolto con la
    libreria \texttt{timezone}.
    \item \textbf{Foto cross-platform}: percorso file su mobile, URL su Web --
    un widget dedicato sceglie come caricarle.
    \item \textbf{Evoluzione dello schema} senza perdere i dati: introdotte le
    \textbf{migrazioni} del database.
    \item \textbf{Coerenza delle statistiche} dopo le eliminazioni: allineamento
    delle cache in memoria.
  \end{itemize}
\end{frame}

\begin{frame}{Demo dell'app}
  \begin{block}{Percorso della demo dal vivo}
    \begin{enumerate}\small
      \item Apertura su \textbf{Parigi} (viaggio in programma) e itinerario.
      \item Scheda \textbf{Valigia}: ``Genera lista'' $\rightarrow$ sacchetti.
      \item Scheda \textbf{Spese}: barra del budget e ritmo di spesa.
      \item \textbf{Timeline} con i tempi liberi.
      \item \textbf{Diario}: aggiunta di una foto a un'attività.
      \item \textbf{Calendario}, \textbf{Ricerca globale} ed \textbf{Emergenza}.
      \item \textbf{Esporta} un viaggio e mostra il menu di condivisione.
    \end{enumerate}
  \end{block}
\end{frame}

% ============================================================
\begin{frame}[plain]
  \begin{center}
    \vspace{1.5cm}
    {\Huge \textcolor{appblue}{\textbf{Grazie!}}}\\[0.6cm]
    {\Large Travel Organizer Offline}\\[0.2cm]
    \small Domande?
  \end{center}
\end{frame}

\end{document}
